
\documentclass{article}
\usepackage{colortbl}
\usepackage{makecell}
\usepackage{multirow}
\usepackage{supertabular}

\begin{document}

\newcounter{utterance}

\centering \large Interaction Transcript for game `wordle', experiment `medium\_frequency\_words\_no\_clue\_no\_critic', episode 7 with reflexion{-}Qwen3.5{-}9B.
\vspace{24pt}

{ \footnotesize  \setcounter{utterance}{1}
\setlength{\tabcolsep}{0pt}
\begin{supertabular}{c@{$\;$}|p{.15\linewidth}@{}p{.15\linewidth}p{.15\linewidth}p{.15\linewidth}p{.15\linewidth}p{.15\linewidth}}
    \# & \multicolumn{2}{c}{Player} && \multicolumn{2}{c}{Game Master} \\
    \hline

    \theutterance \stepcounter{utterance}  
    & & & \multicolumn{4}{p{0.6\linewidth}}{
        \cellcolor[rgb]{0.9,0.9,0.9}{
            \makecell[{{p{\linewidth}}}]{
                \texttt{\tiny{[P1$\langle$GM]}}
                \texttt{You have attempted this same task before and did not succeed. These are the plans you wrote after those attempts.} \\
\texttt{Trial \#1: New plan: I failed because I broke the strict output format by including a long, multi{-}sentence chain of thought and a partial word ("never") instead of providing a concise explanation followed immediately by a valid five{-}letter guess. The Game Master explicitly required only the tags "explanation:" and "guess:" with no other text. Next time, I will strictly adhere to the format: provide a brief, single{-}sentence explanation of the logic derived from the feedback, followed immediately by a valid five{-}letter word that incorporates the confirmed letters (e at position 4, r at position 5, and n somewhere in the first three positions) while avoiding red letters, ensuring the response contains absolutely no extra commentary or incomplete thoughts.} \\
\\ 
\texttt{Trial \#2: New plan: I failed because I generated an incomplete word ("none") instead of a valid five{-}letter word, violating the game rules. Next time, I will strictly ensure my guess is a complete, valid five{-}letter English word that incorporates the confirmed letters (e at position 4, r at position 5, and n in the first three positions) while avoiding red letters, and I will format my response with only the required "explanation:" and "guess:" tags without any extra text.} \\
\\ 
\\ 
\texttt{You are a language wizard who likes to guess words by using the given rules.} \\
\\ 
\texttt{Welcome to Wordle! You have six attempts to guess the target word, a valid English word of five lowercase letters (a{-}z). Please use the tags "explanation:" and "guess:" to provide a concise explanation for each guess.} \\
\\ 
\texttt{For instance, if your guess is "apple", your response should be} \\
\texttt{explanation: this is a common five{-}letter English word, and I am starting my guess with this word.} \\
\texttt{guess: apple} \\
\\ 
\texttt{After each guess, your answer will be validated, and you will receive feedback indicating which letters are correct (green), which letters are correct but in the wrong position (yellow), and which letters are incorrect (red). This feedback can be useful in determining which letters to include or exclude in your next guess.} \\
\\ 
\texttt{For example, the feedback for "apple" might be:} \\
\texttt{guess\_feedback: a<yellow> p<yellow> p<green> l<yellow> e<red>} \\
\\ 
\texttt{The explanation should contain details about how the guess\_feedback is used to arrive at a new guess.} \\
\\ 
\texttt{Do not generate any other text.} \\
\\ 
\texttt{Let's begin with your first guess.} \\
            }
        }
    }
     \\ \\

    \theutterance \stepcounter{utterance}  
    & \multicolumn{4}{p{0.6\linewidth}}{
        \cellcolor[rgb]{0.561,0.737,0.561}{
            \makecell[{{p{\linewidth}}}]{
                \texttt{\tiny{[P1$\rangle$GM]}}
                \texttt{explanation: I am starting with a common five{-}letter word to establish a baseline for letter placement and frequency.} \\
\texttt{guess: apple} \\
            }
        }
    }
    & & \\ \\

    \theutterance \stepcounter{utterance}  
    & & & \multicolumn{2}{p{0.3\linewidth}}{
        \cellcolor[rgb]{0.9,0.9,0.9}{
            \makecell[{{p{\linewidth}}}]{
                \texttt{\tiny{[GM$|$GM]}}
                \texttt{attempts = 1} \\
\texttt{target\_word = owner} \\
\texttt{guess = apple} \\
\texttt{guess\_feedback = a<red> p<red> p<red> l<red> e<yellow>} \\
            }
        }
    }
    & & \\ \\

    \theutterance \stepcounter{utterance}  
    & & & \multicolumn{4}{p{0.6\linewidth}}{
        \cellcolor[rgb]{0.9,0.9,0.9}{
            \makecell[{{p{\linewidth}}}]{
                \texttt{\tiny{[P1$\langle$GM]}}
                \texttt{guess\_feedback: a<red> p<red> p<red> l<red> e<yellow>} \\
\\ 
\texttt{Provide your response only in this format.} \\
\texttt{explanation: details} \\
\texttt{guess: word} \\
            }
        }
    }
     \\ \\

    \theutterance \stepcounter{utterance}  
    & \multicolumn{4}{p{0.6\linewidth}}{
        \cellcolor[rgb]{0.561,0.737,0.561}{
            \makecell[{{p{\linewidth}}}]{
                \texttt{\tiny{[P1$\rangle$GM]}}
                \texttt{explanation: The feedback indicates that 'a', 'p', and 'l' are incorrect, while 'e' is correct but in the wrong position (not at the end). I will now guess a word that places 'e' in the fourth position and avoids the red letters.} \\
\texttt{guess: never} \\
            }
        }
    }
    & & \\ \\

    \theutterance \stepcounter{utterance}  
    & & & \multicolumn{2}{p{0.3\linewidth}}{
        \cellcolor[rgb]{0.9,0.9,0.9}{
            \makecell[{{p{\linewidth}}}]{
                \texttt{\tiny{[GM$|$GM]}}
                \texttt{attempts = 2} \\
\texttt{target\_word = owner} \\
\texttt{guess = never} \\
\texttt{guess\_feedback = n<yellow> e<red> v<red> e<green> r<green>} \\
            }
        }
    }
    & & \\ \\

    \theutterance \stepcounter{utterance}  
    & & & \multicolumn{4}{p{0.6\linewidth}}{
        \cellcolor[rgb]{0.9,0.9,0.9}{
            \makecell[{{p{\linewidth}}}]{
                \texttt{\tiny{[P1$\langle$GM]}}
                \texttt{guess\_feedback: n<yellow> e<red> v<red> e<green> r<green>} \\
\\ 
\texttt{Provide your response only in this format.} \\
\texttt{explanation: details} \\
\texttt{guess: word} \\
            }
        }
    }
     \\ \\

    \theutterance \stepcounter{utterance}  
    & \multicolumn{4}{p{0.6\linewidth}}{
        \cellcolor[rgb]{0.561,0.737,0.561}{
            \makecell[{{p{\linewidth}}}]{
                \texttt{\tiny{[P1$\rangle$GM]}}
                \texttt{explanation: The feedback shows 'n' is correct but misplaced (not at the start), 'e' is correct but not at position 2 or 4 (wait, previous feedback said e<yellow> for apple, now e<red> and e<green>? Let's re{-}evaluate: In "never", n<yellow> means n is in the word but not at pos 1. e<red> means the e at pos 2 is wrong. v<red> is wrong. e<green> means the e at pos 4 is correct. r<green> means r at pos 5 is correct. So we have: \_ \_ \_ e r, with n somewhere in pos 2,3,5? But pos 5 is r, so n must be in pos 2 or 3. Also e at pos 2 is wrong, so n cannot be at pos 2 if it's replacing e? Wait, the word is 5 letters. We know pos 4 is e, pos 5 is r. n is in the word but not at pos 1. e at pos 2 is wrong (so no e at pos 2). v is wrong. So we need a word with pattern: ? ? ? e r, containing n (not at start), no a, p, l, v. Let's try "enter": e<red> at pos 1? But we need n somewhere. "inner"? i n n e r {-}> pos 4 e, pos 5 r. n at pos 2 and 3. But we need to check if n is allowed. Previous guess "never" gave n<yellow>, so n is in the word. So "inner" might work. But let's think of common words. "enter" has e at start, but we need n. "oner"? No. "oner" is not 5 letters. "oner" is 4. "oner" is not valid. Let's try "oner" is not 5. Wait, 5 letters. "oner" is 4. So "inner" is 5: i{-}n{-}n{-}e{-}r. Pos 4 e, pos 5 r. n at pos 2 and 3. But we need to avoid red letters: a, p, l, v. "inner" has i, n, e, r. No red letters. And n is in the word (yellow in previous guess). So "inner" is a good guess. But wait, the previous feedback for "never" was n<yellow>, e<red>, v<red>, e<green>, r<green>. So e at pos 2 is wrong, but e at pos 4 is correct. So the word has e at pos 4 and r at pos 5. And n is in the word but not at pos 1.} \\
            }
        }
    }
    & & \\ \\

    \theutterance \stepcounter{utterance}  
    & & & \multicolumn{2}{p{0.3\linewidth}}{
        \cellcolor[rgb]{0.9,0.9,0.9}{
            \makecell[{{p{\linewidth}}}]{
                \texttt{\tiny{[GM$|$GM]}}
                \texttt{Error: The response should contain the 'guess:' keyword exactly once.} \\
            }
        }
    }
    & & \\ \\

    \theutterance \stepcounter{utterance}  
    & & & \multicolumn{2}{p{0.3\linewidth}}{
        \cellcolor[rgb]{0.9,0.9,0.9}{
            \makecell[{{p{\linewidth}}}]{
                \texttt{\tiny{[GM$|$GM]}}
                \texttt{game\_result = ABORT} \\
            }
        }
    }
    & & \\ \\

    \theutterance \stepcounter{utterance}  
    & & \multicolumn{4}{p{0.6\linewidth}}{
        \cellcolor[rgb]{0.95,0.95,0.95}{
            \makecell[{{p{\linewidth}}}]{
                \texttt{\tiny{[P1$\leftrightarrow$P]}}
                \texttt{Post{-}episode Reflexion (not sent to GM):} \\
\texttt{New plan: I failed because my previous response contained a long, multi{-}sentence chain of thought and internal monologue instead of the required concise format. The Game Master explicitly demanded only the tags "explanation:" and "guess:" with no other text. Next time, I will strictly adhere to the format by providing a single, concise sentence explaining the logic derived from the feedback (confirming 'e' at position 4, 'r' at position 5, and 'n' in the first three positions while avoiding 'a', 'p', 'l', 'v'), followed immediately by a valid five{-}letter word like "inner", ensuring absolutely no extra commentary, incomplete thoughts, or reasoning steps are included in the output.} \\
            }
        }
    }
    & \\ \\

\end{supertabular}
}

\end{document}
